\documentclass{article}
\usepackage{fontspec} 
\usepackage{unicode-math}  % 先加载 unicode-math
\setmathfont{XITS Math}    % 再设置数学字体
%\usepackage[UTF8]{ctex} % 如果需要支持中文

\usepackage{anyfontsize}
\usepackage{xeCJK} % XeLaTeX 推荐 xeCJK，LuaLaTeX 也可用   
\usepackage{setspace}  % 添加这个包
\setstretch{1.2}      % 设置行距为1.5倍

\setCJKmainfont[
    Path = C:/USERS/11252/APPDATA/LOCAL/MICROSOFT/WINDOWS/FONTS/,  % 改用 Windows 字体目录
    UprightFont = {SourceHanSansCN-Light1.otf},  % 使用可变字体
    BoldFont = {SourceHanSansCN-Medium1.otf},
    LetterSpace = 2.0,  % 增加字间距
    WordSpace = 1.2,    % 增加词间距
    %BoldFeatures = {RawFeature={+wght=900}},
    %Scale=1
]{SourceHanSansCN}


\setmainfont[
    Path = C:/USERS/11252/APPDATA/LOCAL/MICROSOFT/WINDOWS/FONTS/,  % 改用 Windows 字体目录
    UprightFont = {SourceHanSansCN-Light1.otf},  % 使用可变字体
    BoldFont = {SourceHanSansCN-Medium1.otf},
    LetterSpace = 2.0,  % 增加字间距
    WordSpace = 1.2,    % 增加词间距
    %BoldFeatures = {RawFeature={+wght=900}},
    %Scale=1
]{SourceHanSansCN}

%\setCJKmainfont[
 %   Path = C:/USERS/11252/APPDATA/LOCAL/MICROSOFT/WINDOWS/FONTS/,
 %   UprightFont = {SourceHanSerifCN-Light-5.otf},
 %   BoldFont = {SourceHanSerifCN-Medium-6.otf},
%    Scale=1
%]{SourceHanSerifCN}

%\setmainfont{SourceHanSerif}  % 设置英文字体

% 处理冲突的命令
\let\oldbackepsilon\backepsilon
\let\backepsilon\relax
\let\oldeth\eth
\let\eth\relax
\let\olddigamma\digamma
\let\digamma\relax
\usepackage{float}

\usepackage[a4paper, left=0.1in, right=0.1in, top=0.05in, bottom=0.05in]{geometry} % 设置四边页边距为0.1英寸{geometry} % 设置页边距为0.5英寸
\usepackage{enumitem} % 用于自定义列表
\usepackage{amsmath}  % 数学公式支持
\usepackage{tikz}
\usetikzlibrary{arrows.meta, positioning}
\setlength{\parindent}{0pt} % 不缩进
\usepackage{etoolbox} % 用于列表操作
%调整类代码
\usepackage{comment}
\usepackage{tocloft} % 用于定制目录 样式
\usepackage{hyperref} %不需要目录盒子注释这
\usepackage{ifthen}
\usepackage{tikz}
\usetikzlibrary{arrows.meta}
\usepackage{titletoc}
\usepackage{graphicx}
\usepackage{float}

\usepackage{amssymb}  % 提供数学符号，包括\therefore
%目录设定
%快捷键 CMD + / 注释
%  \renewcommand{\cftdot}{} % 隐藏目录中的页码
%  \cftpagenumbersoff{section}
%  \cftpagenumbersoff{subsection}
%  \makeatletter
%  \renewcommand{\@pnumwidth}{0em}  % 设置页码宽度为0
%  \renewcommand{\@tocrmarg}{0em}   % 设置右边距为0
%  \makeatother
 

\usepackage{environ}

\newif\ifshowcontent  % 定义一个条件判断变量
\showcontenttrue    % 设置为 false，隐藏所有 question 环境的内容

\newcounter{question} % 题目计数器
\setcounter{question}{0}
\NewEnviron{question}{%
    \ifshowcontent
        \par\stepcounter{question}\textbf{\thequestion.} \BODY\par % 如果显示内容，则输出
    \fi
}

\NewEnviron{choices}{%
    \ifshowcontent
        \begin{enumerate}[label=\Alph*., leftmargin=*]
            \BODY
        \end{enumerate}\par
    \fi
}

\NewEnviron{correctanswer}{%
    \ifshowcontent
        \textbf{答案:}\BODY\par
    \fi
}



\newenvironment{explanation}
{\textbf{解析:}\par}
{\par}



% 新增考察方面命令
\newcommand{\checklist}[1]{
    \textbf{考察:} #1 \par % 在题目下方显示考察方面
    \addcontentsline{toc}{subsection}{题号\thequestion 考察: #1} % 将考察内容添加到目录
    \vspace{2em} % 添加1em的垂直空间
}

\graphicspath{{images/}} %统一


\begin{document}
 %\fontsize{22}{31}\selectfont
 \tableofcontents % 添加目录
\newpage
%\pagestyle{empty}




\section{考点1:预警指标}
\setcounter{question}{0}

\begin{question}
    Firms generally use a stoplight system in representing and communicating their performance against the thresholds of their EWIs. A green indicator means the measure is within normal bounds; An amber indicator means the measure should be investigated further; A red indicator should be a source for significant concern and may warrant an immediate response. Which of the following is the most suitable way to start a calibration of these thresholds?
\end{question}

\begin{choices}
    \item The boundaries for which an EWI moves from green to yellow should be narrow enough
    \item Historical data can inform the calibration of thresholds
    \item Practitioners observe that this stoplight system is effective because it can be set very wide
    \item The threshold should not be subject to back-testing
\end{choices}

\begin{correctanswer}
    B
\end{correctanswer}

\begin{explanation}
    \textbf{A选项错误。}
    这个说法不准确：
    \begin{itemize}
        \item 从绿色到黄色的边界不应过窄
        \item 过窄会导致过度预警
    \end{itemize}

    \textbf{B选项正确。}
    这个说法正确：
    \begin{itemize}
        \item 历史数据可以指导阈值校准
        \item 这是最合适的起点
        \item 有助于设定合理的阈值
    \end{itemize}

    \textbf{C选项错误。}
    这个说法不准确：
    \begin{itemize}
        \item 信号灯系统不应设置过宽
        \item 过宽会降低预警效果
    \end{itemize}

    \textbf{D选项错误。}
    这个说法不准确：
    \begin{itemize}
        \item 阈值应该进行回测
        \item 回测有助于验证阈值有效性
    \end{itemize}
\end{explanation}

\checklist{预警指标阈值校准}


\begin{question}
    The framework of Early Warning Indicators can be summarized as M.E.R.I.T" and "M" means "Measures". Which of the following is true about a good Early Warning Indicator (EWI) measures?
\end{question}

\begin{choices}
    \item Forward-looking and sharp measures
    \item Measures are not linked from the escalation process
    \item Leading indicators provide Information and signal potential stress after the occurrence of an actual event
    \item Banks should consider only internal but not external measures
\end{choices}

\begin{correctanswer}
    A
\end{correctanswer}

\begin{explanation}
    \textbf{A选项正确。}
    这个说法正确：
    \begin{itemize}
        \item 好的预警指标应具有前瞻性
        \item 指标应具有敏锐性
        \item 能及时反映潜在风险
    \end{itemize}

    \textbf{B选项错误。}
    这个说法不准确：
    \begin{itemize}
        \item 指标应与升级流程相关联
        \item 升级流程是预警指标的重要组成部分
    \end{itemize}

    \textbf{C选项错误。}
    这个说法不准确：
    \begin{itemize}
        \item 领先指标应在事件发生前提供信息
        \item 而不是在事件发生后
    \end{itemize}

    \textbf{D选项错误。}
    这个说法不准确：
    \begin{itemize}
        \item 银行应同时考虑内部和外部指标
        \item 不应仅关注内部指标
    \end{itemize}
\end{explanation}

\checklist{预警指标特征分析}


\begin{question}
    Venkat explains that firms often use a stoplight system to manage their thresholds: "Firms generally use a stoplight system in representing and communicating their performance against the thresholds of their EWIs. A green indicator means that the measure is within normal bounds. A measure that is classified as amber according to the threshold framework should be investigated further while a red indicator should be a source for significant concern and may warrant an immediate response." Which of the following is the BEST way to start the exercise of calibration of these thresholds?
\end{question}

\begin{choices}
    \item Thresholds are ultimately subjective
    \item Historical data can inform the calibration of thresholds
    \item Practitioners observe that this stoplight system is obsolete and ineffective such that thresholds are moot
    \item The firm should rely on the specific regulatory instructions, in particular BCBS 2008, for calibration of the thresholds
\end{choices}

\begin{correctanswer}
    B
\end{correctanswer}

\begin{explanation}
    \textbf{A选项错误。}
    这个说法不准确：
    \begin{itemize}
        \item 阈值不应完全主观
        \item 需要客观依据支持
    \end{itemize}

    \textbf{B选项正确。}
    这个说法正确：
    \begin{itemize}
        \item 历史数据可以指导阈值校准
        \item 这是最合适的起点
        \item 有助于设定合理的阈值
    \end{itemize}

    \textbf{C选项错误。}
    这个说法不准确：
    \begin{itemize}
        \item 信号灯系统并非过时无效
        \item 阈值仍然很重要
    \end{itemize}

    \textbf{D选项错误。}
    这个说法不准确：
    \begin{itemize}
        \item 不应仅依赖监管指引
        \item 需要结合实际情况
    \end{itemize}
\end{explanation}

\checklist{预警指标阈值校准}




\newpage



\section{考点2:日内流动性风险管理}
\setcounter{question}{0}


\begin{question}
    Which of the following factors is least likely to increase a bank treasurer's challenge in managing cash in order to stay within its intraday credit limits?
\end{question}

\begin{choices}
    \item Cash flow volatility
    \item Credit quality of assets
    \item Insufficient real-time data
    \item Securities price volatility
\end{choices}

\begin{correctanswer}
    B
\end{correctanswer}

\begin{explanation}
    \textbf{A选项错误。}
    这个说法不准确：
    \begin{itemize}
        \item 现金流波动会增加管理难度
        \item 需要更频繁地调整头寸
    \end{itemize}

    \textbf{B选项正确。}
    这个说法正确：
    \begin{itemize}
        \item 银行资产信用质量很高
        \item 包括现金余额和流动资产
        \item 这些资产几乎没有违约风险
        \item 不会增加日间流动性管理难度
    \end{itemize}

    \textbf{C选项错误。}
    这个说法不准确：
    \begin{itemize}
        \item 实时数据不足会增加管理难度
        \item 影响决策的及时性
    \end{itemize}

    \textbf{D选项错误。}
    这个说法不准确：
    \begin{itemize}
        \item 证券价格波动会增加管理难度
        \item 影响资产估值和流动性
    \end{itemize}
\end{explanation}

\checklist{日间流动性管理分析}


\begin{question}
    In the context of characteristics of an effective governance structure in controlling intraday liquidity within a bank, there is emphasis on expertise in which line of defense?
\end{question}

\begin{choices}
    \item Treasury
    \item Internal audit
    \item Information technology
    \item Corporate risk management
\end{choices}

\begin{correctanswer}
    D
\end{correctanswer}

\begin{explanation}
    \textbf{A选项错误。}
    这个说法不准确：
    \begin{itemize}
        \item 资金部是第一道防线
        \item 负责日常运营管理
    \end{itemize}

    \textbf{B选项错误。}
    这个说法不准确：
    \begin{itemize}
        \item 内部审计是第三道防线
        \item 负责独立监督
    \end{itemize}

    \textbf{C选项错误。}
    这个说法不准确：
    \begin{itemize}
        \item 信息技术不是防线
        \item 是支持性功能
    \end{itemize}

    \textbf{D选项正确。}
    这个说法正确：
    \begin{itemize}
        \item 企业风险管理是第二道防线
        \item 强调专业知识和技能
        \item 负责风险监控和管理
    \end{itemize}
\end{explanation}

\checklist{银行治理结构分析}


\begin{question}
    All risk management frameworks start with a governance structure that defines the roles and responsibilities of various bank employees and committees in overseeing risk-related activities. Effective governance includes the oversight of intraday liquidity risk. Which of the following statements is true about intraday liquidity risk management?
\end{question}

\begin{choices}
    \item Payments on large value payment systems should not be considered as a intraday liquidity
    \item Roles and responsibilities should be defined by Treasury department
    \item PCS Systems should only be a source of funds; if PCS becomes a use of funds, then a yellow flag should be triggered
    \item Intraday liquidity risk should be incorporated in the risk taxonomy and is a component of risk self-assessment including settlement risks
\end{choices}

\begin{correctanswer}
    D
\end{correctanswer}

\begin{explanation}
    \textbf{A选项错误。}
    这个说法不准确：
    \begin{itemize}
        \item 大额支付系统是重要的日间信用使用
        \item 应纳入日间流动性管理
    \end{itemize}

    \textbf{B选项错误。}
    这个说法不准确：
    \begin{itemize}
        \item 角色和职责应由资产负债委员会(ALCO)定义
        \item 不是由资金部定义
    \end{itemize}

    \textbf{C选项错误。}
    这个说法不准确：
    \begin{itemize}
        \item 支付、清算和结算(PCS)系统可以是资金来源或使用
        \item 不应仅作为资金来源
    \end{itemize}

    \textbf{D选项正确。}
    这个说法正确：
    \begin{itemize}
        \item 日间流动性风险应纳入风险分类
        \item 是风险自评估的组成部分
        \item 包括结算风险
    \end{itemize}
\end{explanation}

\checklist{日间流动性风险管理}


\begin{question}
    Each of the following is a measure for quantifying and/or monitoring risk levels except which is a measure for understanding intraday flows?
\end{question}

\begin{choices}
    \item Total payments
    \item Client intraday credit usage
    \item Intraday credit relative to tier 1 capital
    \item Daily maximum intraday liquidity usage
\end{choices}

\begin{correctanswer}
    A
\end{correctanswer}

\begin{explanation}
    \textbf{A选项正确。}
    这个说法正确：
    \begin{itemize}
        \item 总支付是用于理解日间流动的指标
        \item 不是风险量化或监控指标
    \end{itemize}

    \textbf{B选项错误。}
    这个说法不准确：
    \begin{itemize}
        \item 客户日间信用使用是风险监控指标
        \item 用于量化风险水平
    \end{itemize}

    \textbf{C选项错误。}
    这个说法不准确：
    \begin{itemize}
        \item 日间信用相对于一级资本的比率是风险监控指标
        \item 用于量化风险水平
    \end{itemize}

    \textbf{D选项错误。}
    这个说法不准确：
    \begin{itemize}
        \item 每日最大日间流动性使用是风险监控指标
        \item 用于量化风险水平
    \end{itemize}
\end{explanation}

\checklist{日间流动性指标分析}

\begin{question}
    Which of the following is a USE of intraday liquidity?
\end{question}

\begin{choices}
    \item Income funds flow
    \item Term repo (as the repo seller)
    \item Funding of nostro accounts (at correspondent bank outside home market)
    \item Intraday credit (Federal Reserve unsecured committed line of credit, LOC)
\end{choices}

\begin{correctanswer}
    C
\end{correctanswer}

\begin{explanation}
    \textbf{A选项错误。}
    这个说法不准确：
    \begin{itemize}
        \item 收入资金流是日间流动性的来源
        \item 不是使用
    \end{itemize}

    \textbf{B选项错误。}
    这个说法不准确：
    \begin{itemize}
        \item 定期回购(作为回购卖方)是日间流动性的来源
        \item 不是使用
    \end{itemize}

    \textbf{C选项正确。}
    这个说法正确：
    \begin{itemize}
        \item 代理行账户资金(在母国市场外的代理行)是日间流动性的使用
        \item 需要动用日间流动性
    \end{itemize}

    \textbf{D选项错误。}
    这个说法不准确：
    \begin{itemize}
        \item 日间信用(美联储无担保承诺信用额度)是日间流动性的来源
        \item 不是使用
    \end{itemize}
\end{explanation}

\checklist{日间流动性使用分析}



\newpage


\section{考点3:流动性风险报告}
\setcounter{question}{0}
\begin{question}
    Which of the following is the most likely output of a bank's liquidity stress test report?
\end{question}

\begin{choices}
    \item Cash flow survival horizon
    \item Net interest income projection
    \item Leverage ratio of unsecured debt to equity
    \item Return on equity versus threshold and target
\end{choices}

\begin{correctanswer}
    A
\end{correctanswer}

\begin{explanation}
    \textbf{A选项正确。}
    这个说法正确：
    \begin{itemize}
        \item 现金流生存期是流动性压力测试报告最重要的输出
        \item 反映银行在压力情况下的生存能力
        \item 是流动性风险管理的核心指标
    \end{itemize}

    \textbf{B选项错误。}
    这个说法不准确：
    \begin{itemize}
        \item 净利息收入预测不是压力测试的主要输出
        \item 属于盈利性分析
    \end{itemize}

    \textbf{C选项错误。}
    这个说法不准确：
    \begin{itemize}
        \item 无担保债务与权益的杠杆比率不是压力测试的主要输出
        \item 属于资本结构分析
    \end{itemize}

    \textbf{D选项错误。}
    这个说法不准确：
    \begin{itemize}
        \item 股本回报率与阈值和目标的比较不是压力测试的主要输出
        \item 属于盈利性分析
    \end{itemize}
\end{explanation}

\checklist{流动性压力测试分析}





\newpage



\section{考点4:流动性压力测试}
\setcounter{question}{0}

\begin{question}
    In the context of deposit outflows, which behavioral assessment factor is relevant for individual, small business, and commercial/institutional customers of the bank?
\end{question}

\begin{choices}
    \item Credit usage
    \item FDIC coverage
    \item Industry segment
    \item Relationship tenure
\end{choices}

\begin{correctanswer}
    D
\end{correctanswer}

\begin{explanation}
    \textbf{A选项错误。}
    这个说法不准确：
    \begin{itemize}
        \item 信用使用主要适用于小型企业和商业/机构客户
        \item 不适用于个人客户
    \end{itemize}

    \textbf{B选项错误。}
    这个说法不准确：
    \begin{itemize}
        \item FDIC保险仅适用于个人和小型企业
        \item 不适用于商业/机构客户
    \end{itemize}

    \textbf{C选项错误。}
    这个说法不准确：
    \begin{itemize}
        \item 行业细分不适用于个人客户
        \item 主要适用于商业/机构客户
    \end{itemize}

    \textbf{D选项正确。}
    这个说法正确：
    \begin{itemize}
        \item 关系权属适用于所有三类客户
        \item 是共同的行为评估因素
        \item 反映客户与银行的长期关系
    \end{itemize}
\end{explanation}

\checklist{存款流出行为分析}


\begin{question}
    Which of the following types of liquidity is available in the form of the institution's liquidity asset buffer and represents liquidity available to meet general financial obligations under a stress scenario?
\end{question}

\begin{choices}
    \item Operational
    \item Strategic
    \item Contingent
    \item Restricted
\end{choices}

\begin{correctanswer}
    C
\end{correctanswer}

\begin{explanation}
    \textbf{A选项错误。}
    这个说法不准确：
    \begin{itemize}
        \item 运营流动性不可用于满足一般财务义务
        \item 在压力测试中不可用
    \end{itemize}

    \textbf{B选项错误。}
    这个说法不准确：
    \begin{itemize}
        \item 战略流动性不可用于满足一般财务义务
        \item 在压力测试中不可用
    \end{itemize}

    \textbf{C选项正确。}
    这个说法正确：
    \begin{itemize}
        \item 或有流动性可用于满足一般财务义务
        \item 以机构流动性资产缓冲形式存在
        \item 在压力情景下可用
    \end{itemize}

    \textbf{D选项错误。}
    这个说法不准确：
    \begin{itemize}
        \item 限制性流动性不可用于满足一般财务义务
        \item 在压力测试中不可用
    \end{itemize}
\end{explanation}

\checklist{流动性类型分析}




\newpage



\section{考点5:应急资金规划}
\setcounter{question}{0}


\begin{question}
    In the context of governance and oversight of a contingency funding plan (CFP), which role has direct oversight of the liquidity crisis team (LCT)?
\end{question}

\begin{choices}
    \item Board of directors
    \item Risk management
    \item Corporate treasury
    \item Management committee
\end{choices}

\begin{correctanswer}
    D
\end{correctanswer}

\begin{explanation}
    \textbf{A选项错误。}
    这个说法不准确：
    \begin{itemize}
        \item 董事会不直接监督LCT
        \item 仅提供建议和协助
    \end{itemize}

    \textbf{B选项错误。}
    这个说法不准确：
    \begin{itemize}
        \item 风险管理不直接监督LCT
        \item 负责风险监控
    \end{itemize}

    \textbf{C选项错误。}
    这个说法不准确：
    \begin{itemize}
        \item 公司资金部不直接监督LCT
        \item 负责日常资金管理
    \end{itemize}

    \textbf{D选项正确。}
    这个说法正确：
    \begin{itemize}
        \item 管理委员会(包括高级管理层)负责直接监督LCT
        \item 是应急融资计划治理的核心
        \item 确保LCT有效运作
    \end{itemize}
\end{explanation}

\checklist{应急融资计划治理}


\begin{question}
    Which of the following items is an example of a contingent action that could be taken by a bank during a stress situation?
\end{question}

\begin{choices}
    \item Securitizing assets
    \item Decreasing lending rates
    \item Increasing capital distributions
    \item Shifting from longer-term to shorter-term funding sources
\end{choices}

\begin{correctanswer}
    A
\end{correctanswer}

\begin{explanation}
    \textbf{A选项正确。}
    这个说法正确：
    \begin{itemize}
        \item 资产证券化是银行的现金来源
        \item 增加流动性
        \item 是压力情况下的或有行动
    \end{itemize}

    \textbf{B选项错误。}
    这个说法不准确：
    \begin{itemize}
        \item 降低贷款利率会鼓励增加放贷
        \item 导致现金流出
        \item 减少流动性
    \end{itemize}

    \textbf{C选项错误。}
    这个说法不准确：
    \begin{itemize}
        \item 增加资本分配会减少流动性
        \item 不是压力情况下的适当行动
    \end{itemize}

    \textbf{D选项错误。}
    这个说法不准确：
    \begin{itemize}
        \item 从长期转向短期融资不是或有行动
        \item 正确的方向是从短期转向长期融资
    \end{itemize}
\end{explanation}

\checklist{压力情况应对措施}


\begin{question}
    Which level of escalation within a contingency funding plan (CFP) has more emphasis on analyzing the causes of liquidity deterioration closely?
\end{question}

\begin{choices}
    \item Level 1
    \item Level 2
    \item Level 3
    \item Level 4
\end{choices}

\begin{correctanswer}
    B
\end{correctanswer}

\begin{explanation}
    \textbf{A选项错误。}
    这个说法不准确：
    \begin{itemize}
        \item 1级是初始预警级别
        \item 不需要详细分析原因
    \end{itemize}

    \textbf{B选项正确。}
    这个说法正确：
    \begin{itemize}
        \item 2级是更高级别的预警
        \item 系统性或特例事件明显影响业务和流动性
        \item 需要详细分析当前流动性和恶化原因
    \end{itemize}

    \textbf{C选项错误。}
    这个说法不准确：
    \begin{itemize}
        \item 3级是严重预警级别
        \item 重点在于采取行动而非分析原因
    \end{itemize}

    \textbf{D选项错误。}
    这个说法不准确：
    \begin{itemize}
        \item 4级是最高预警级别
        \item 重点在于危机管理而非分析原因
    \end{itemize}
\end{explanation}

\checklist{应急融资计划升级}


\begin{question}
    What is a key design consideration in developing a contingency funding plan (CFP)?
\end{question}

\begin{choices}
    \item Aligned to business profile
    \item Supported by a backup plan
    \item Inclusive of all shareholder groups
    \item Can be used in both normal and stressed states
\end{choices}

\begin{correctanswer}
    A
\end{correctanswer}

\begin{explanation}
    \textbf{A选项正确。}
    这个说法正确：
    \begin{itemize}
        \item CFP应与公司的业务和风险状况相匹配
        \item 包括业务活动、产品和资产类别
        \item 是制定CFP的关键考虑因素
    \end{itemize}

    \textbf{B选项错误。}
    这个说法不准确：
    \begin{itemize}
        \item 备用计划不是CFP的关键设计考虑因素
        \item CFP本身就是应对压力的计划
    \end{itemize}

    \textbf{C选项错误。}
    这个说法不准确：
    \begin{itemize}
        \item CFP应包含适当的利益相关者群体
        \item 不仅限于股东
        \item 还包括管理委员会、业务部门等
    \end{itemize}

    \textbf{D选项错误。}
    这个说法不准确：
    \begin{itemize}
        \item CFP本质上仅用于压力状态
        \item 不适用于正常状态
    \end{itemize}
\end{explanation}

\checklist{应急融资计划设计}


\begin{question}
    Alan Chen, a staff of JY bank, is considering developing an appropriate contingency funding plan, all of the following statements should be considered in building Contingency Funding Plan EXCEPT:
\end{question}

\begin{choices}
    \item Aligned to business and risk profiles
    \item Integrated with broader risk management frameworks
    \item Theoretical basis for precoded instruction set
    \item Inclusive of appropriate stakeholder groups
\end{choices}

\begin{correctanswer}
    C
\end{correctanswer}

\begin{explanation}
    \textbf{A选项错误。}
    这个说法不准确：
    \begin{itemize}
        \item 与业务和风险状况相匹配是CFP的正确设计考虑因素
        \item 是五个关键考虑因素之一
    \end{itemize}

    \textbf{B选项错误。}
    这个说法不准确：
    \begin{itemize}
        \item 与更广泛的风险管理框架整合是CFP的正确设计考虑因素
        \item 是五个关键考虑因素之一
    \end{itemize}

    \textbf{C选项正确。}
    这个说法正确：
    \begin{itemize}
        \item 预编码指令集的理论基础不是CFP的设计考虑因素
        \item 五个正确的设计考虑因素是：
        \begin{enumerate}
            \item 与业务和风险状况相匹配
            \item 与更广泛的风险管理框架整合
            \item 可操作且灵活的行动手册
            \item 包含适当的利益相关者群体
            \item 有沟通计划支持
        \end{enumerate}
    \end{itemize}

    \textbf{D选项错误。}
    这个说法不准确：
    \begin{itemize}
        \item 包含适当的利益相关者群体是CFP的正确设计考虑因素
        \item 是五个关键考虑因素之一
    \end{itemize}
\end{explanation}

\checklist{应急融资计划设计考虑因素}




\newpage





\end{document}